% ============================================================
%  G. Brandes, "Dyret i Mennesket"
%  Af Dagens Krønike, Maj–Juni 1890, pp. 494–551
%  (Reprinted in: Samlede Skrifter, vol. VII)
%  TRANSLATION (English)
% ============================================================
\documentclass[12pt,a4paper]{article}
\usepackage[utf8]{inputenc}
\usepackage[T1]{fontenc}
\usepackage[english]{babel}
\usepackage{microtype}
\usepackage{setspace}
\usepackage[margin=1.2in]{geometry}
\usepackage{fancyhdr}
\usepackage{hyperref}

\hypersetup{
  pdftitle={The Animal in the Human Being},
  pdfauthor={G. Brandes},
  hidelinks
}

\pagestyle{fancy}
\fancyhf{}
\rhead{Brandes, \emph{The Animal in the Human Being} (1890)}
\cfoot{\thepage}

\setstretch{1.3}

\title{\textbf{The Animal in the Human Being}}
\author{G.\ Brandes\\[4pt]
  \small Translated from the Danish}
\date{1890}

\begin{document}
\maketitle
\thispagestyle{fancy}

\begin{center}
  \small (Bourget: \textit{The Disciple}. Maupassant: \textit{Useless Beauty}.
  Tolstoy: \textit{The Kreutzer Sonata}. Zola: \textit{The Beast in Man}.)
\end{center}

\bigskip

It is striking to what a high degree the literature of recent times,
in its most distinguished European representatives, is preoccupied
with the idea of a duality in our nature.

In the work of a single French author, Paul Bourget, this duality in
human nature appears with a scientific grounding. For Bourget is, even
as a creative writer, the soul-probing critic he began by being. In
explaining the instinctual and emotional lives of his characters he
never proceeds in a purely artistic manner; he always points with the
investigator's probe.

If one surveys what he has produced as a writer, from his first novella
\textit{L'irréparable} to his latest novel \textit{Le disciple}, there
are accordingly certain striking common features. Everything turns, at
bottom, on the doubling of the self. His first novella is the story of
a peculiar young woman whose life's critical event is a nocturnal
assault to which she has been subjected in the loneliness of night by
an unscrupulous man with whom she has engaged in an innocent flirtation.
The tale begins, very characteristically, with a conversation between
the author and a professor of psychology who has written a work on the
dissolution of the associations of ideas.

Now, since what we call our self---this idea of something inner and
permanent that is united with a body---is an idea that arises through
the linking of ideas, it is clear that what the professor has in
reality studied are the forms of the self's dissolution, of its
duality, as this appears in the numerous diseases of will and mind. In
the last novel Bourget has published, we encounter once again the
professor of psychology, and once again the same subject. The work that
the young criminal Robert Greslou, accused of murder, has submitted to
a celebrated professor in order to learn his opinion, is a study of the
duality of the self.

In the novella it is first developed wherein the self's simplest
doubling consists. The self is partly conscious, partly unconscious.
This explains in general the mistakes and missteps of life. There is
hidden within us a creature we do not know, and of which we never know
whether it is not the very opposite of the being we believe ourselves
to be. On this rest the remarkable reversals of mood and conduct that
we experience or observe. One person works toward a goal upon which
he imagines his happiness to depend, and when it is reached he
discovers that he has misjudged the secret and true demands of his
emotional life. Another goes about looking like a calm person, and
acts accordingly; suddenly the image of his misery rises up before
him---a persistent sorrow or persistent danger that he has sought to
forget, but which can overwhelm him and give his life a different
direction.

In the novel, the main character furnishes a complete investigation of
his own being, tracing every state of his mind back to a scientific
schema: heredity, the influence of environment, transplantation of
the person, and so on. He determines his dominant characteristic as
the capacity and urge for self-doubling: there are always in him, as
it were, two clearly distinct personalities, one who comes and goes,
acts and feels, and one who, as spectator, observes the other's
conduct of life. He is unable to say which of these two persons is
his true self. He tries to derive this duality in himself from a
mixture of races, from descent from a border people. From childhood he
has taken pleasure in pretence and untruth---not in order to boast,
but simply in order to be, in his imagination and in the consciousness
of his surroundings, someone else. He takes the greatest pleasure in
expressing views that are the complete opposite of those he regards
as true, and for the same strange reason. His delight is to play a
role alongside his true nature. His whole mendacity of bearing, his
joy in deceiving, seducing, outwitting, and appropriating a young
noble-minded woman for whom he feels no tenderness---all of it is a
consequence of the unconquerable duality in the original disposition
of his human nature.

As one can see, the bifurcation that preoccupies the scientifically
formed Bourget is one that can be fully understood only through
philosophical means. It admittedly corresponds, to some degree, to the
split in our nature that in former times was found between what was
called the corporeal and what was called the spiritual, but it does
not coincide with that split in any way. It is conceived in the
old-fashioned manner only insofar as Bourget, in his most recent
writings, seems to wish to glorify religious devotion as the only
remedy against it.

\section*{I}

Other significant authors of the age who likewise dwell on the idea
of a duality in our nature have given this bifurcation an expression
that belongs more properly to the Middle Ages. They speak not
infrequently of the animal in the human being and of the human in the
human being as of two different worlds coupled together within us.
They show themselves occupied with this relationship in much the same
way as the theologians were with the relationship between the divine
and the human nature in the God-man. For quite a few contemporary
authors, the human creature is a kind of centaur---half a higher
being, half an animal---which, when it loses itself in soaring
thoughts, is awakened from its dream by the sound of its own hooves.

Those who brood over the animal in the human being will either merely
depict it, paint it in its power and repulsiveness, or they will
combat it---sometimes even seeking to eradicate it---without ever
managing sharply to indicate the dividing line where the animal ends
and the human begins. For there are only too many activities and
drives that the human being shares with the lower creatures.

Thought-provoking, now, is the fundamental disagreement among our
age's most significant authors regarding the question of whether the
so-called animal element is an unconquered remnant of the original
state of nature, of the savage era and the days of barbarism, or
whether it, as we know it, is a product of culture, nurtured by
over-civilization, so that we can combat it only by, as the saying
goes, returning to nature---that is, to simpler, more rural
conditions.

A writer like Maupassant or Zola inclines toward the first view; for
them the animal element is the original savagery. A writer like
Strindberg or Tolstoy embraces the second perspective; for them the
animal element is the unnaturalness of our false civilization.

Tolstoy's new story \textit{The Kreutzer Sonata}---banned in Russia
but widely read---is a purely moralizing work, related in this respect
to his edifying legends and stories for the Russian common people. Only
the book addresses a readership that belongs to the most enlightened
in the world.

Here as in \textit{Ghosts} and \textit{A Gauntlet}, sexual purity and
impurity is what everything turns on; but there is in the great
Russian's perspective a wilder, more passionately consistent
thoroughness than in the Norwegians. Ibsen in \textit{Ghosts} had
painted terrifying consequences of male licentiousness for the home
and the next generation; Bjørnson had let a young woman make demands
for the man's unconditional abstinence before marriage; it seemed
impossible to go higher or further in one's demands in this field.
But Tolstoy, like a true Russian, runs the line out to its end---this
is the distinguishing mark by which, in my book \textit{Impressions of
Russia}, I tried to paint Russian distinctiveness. He takes up again
the demand for the man's purity before marriage, but augments it with
a demand for so-called purity in marriage that is new.

To be sure, Tolstoy does not speak in his own name in the new story.
But partly through a conversation in a railway carriage, in which the
views that emerge victorious seem likely to lie close to the author's
own, and partly through the main character's long account of his life,
which takes up the greater part of the work, Tolstoy's sympathy with
the outlook on life that finds expression can be traced. It is not
probable that he would wish to maintain any essential disagreement
between Pozdnyshev's fundamental perspective and his own. In any case
he has done nothing to give the reader the impression that such a
disagreement obtains. All that can be said is that in his writing
\textit{What I Believe} he has admittedly moved along the same paths
as his main character, but has not gone quite as far.

Tolstoy passionately upholds the old Christian conception of the
virginal state as the genuinely high and valuable one. He has forceful
words about the foolishness of a social spirit that throws many a young
lovable girl into the arms of an impure and depraved man merely so
that she shall not suffer the supposed shame of remaining a virgin---
which in Tolstoy's eyes is an honour. This alone is radical enough in
feeling, even if the matter has not been thought through.

Marriage or non-marriage, a ceremony or no ceremony, is for Tolstoy
a matter of complete indifference; he places no weight on the form.
But the first woman a man has bound himself to, the first man a woman
has given her trust to---to him and her they are bound for life,
indissolubly. They are united in a relation of duty. Whether there is
joy in it is not asked; no weight is placed upon it.

The modern view that it is love that grounds marriage and gives it
worth strikes Tolstoy as a frivolity---much as it appeared to people
of the past, and indeed at bottom still to Hegel with his sober
disposition. Tolstoy seems to have arrived at the conviction that,
in the opinion of various dark-seeing people, is the correct one:
that what was previously called marriage, and is still mostly called
that, is only possible where everything is built on authority and
obedience, where the parents choose a husband for their daughter and
obtain a wife for their son when the time for it has come, and where
the wife trembles before the husband it is her duty to love. For this
is of course the classical form of marriage: no choice, no divorce;
one master, one will, corporal punishment in cases of insubordination;
much domestic discipline. If one is unwilling to accept this, he seems
to think, one gradually arrives at making love sovereign, making
sexual union a private matter, and limiting social activity to
securing, to the best of one's ability, the welfare of children.

The most conservative but not the most adequate interpretation of
\textit{The Kreutzer Sonata} is that Tolstoy wanted to say: ``As men
are now, they are made unfit for married life by their conduct in
youth.'' It is undeniable that he, quite like Bjørnson, is very much
preoccupied with the wretchedness and degradation of the life that men
are said generally to lead before entering into marriage. Tolstoy has
even in his own confessional writings given accounts of a rather
disturbing nature which, even if they exaggerate, show that he knows
from his own experience what he judges most harshly:

\begin{quote}
``It is impossible for me to recall those years without being overcome
by a feeling of horror, disgust, and pain. There is no vice to which
I could not have abandoned myself at that time, no crime I would not
have been capable of committing. Lying, theft, debaucheries of every
kind, violence, and murder---these are things I have made myself
guilty of in all their forms\ldots\ On the estate where I lived I
squandered in revelry and card-playing everything my serfs had earned
for me by their labour---and in addition I tormented and punished them
on every occasion, sacrificed them to my excesses, deceived them, sold
them, and so on.''
\end{quote}

Compare with this, in \textit{The Kreutzer Sonata}, Pozdnyshev's
description of how he ``wallowed in the filth of debauchery.'' The
kernel here is the same as in Bjørnson's accounts of a related kind:
men---all men of the better-off classes at least---live in their youth
a downright swinish life, crude, loveless, with women indiscriminately.

The certain thing about the matter is that accounts and statements of
this kind can in truth have a lively interest only for those whose
development has not run ahead of these authors'. A not inconsiderable
number of male readers' impression will be this: we do not feel
addressed. Many a person who is not suspected of being a hypocrite will
say to himself quietly: I am tired of this talk. It is talk for the
spiritual common people. However excellently such books may be written,
in this respect they are books for the common people, which merely
impress upon one anew that all genuine art addresses itself to the most
developed of the age.

We see in different countries related literary phenomena of this kind.
There arises the younger Dumas who, by his own account, by no means
lived like a saint in his youth ``unless one were to think of Saint
Augustine in the first phase of his life''; he sets his life in order
and, on the threshold of old age, preaches the strictest sexual
morality---spicing it, from old habit, with bold figures of speech and
risqué allusions. There arises Bjørnson who, without having Dumas's
Parisian experiences of youth, has developed into an equally stubborn
teacher of duty and who has it from a genuine Norwegian that he is
more fresh and sincere than actually refined. He proclaims to us the
sexual decline of men. He seems, it appears, to have seen very much
crudeness around him, and if he can diminish it, that is meritorious.
Whether it is diminished by a touring moral agitation of the kind he
embarked upon in his time is an open question. What is certain is
that he has diminished the artistic value of his great didactic novels
by the pastoral care that dominates them.

Now comes Tolstoy. He has behind him a Russian ``junker's'' civil and
military experiences, and draws conclusions from these for men who
have never been ``junkers'', never led a junker's life, and perhaps
never displayed crudeness in relation to the other sex. But upon such
readers he, like the other preachers of morality in their struggle
against social morals, necessarily makes the impression of Puritans
attacking Papists. The developed reader has no interest in internal
theological feuds. Monkish squabbling!---as Ulrich von Hutten wrote
about Luther's first attack on the Pope.

There surely exist even in our day (or already in our day) men who
feel, on decisive points, like Greeks from the best period of ancient
Hellas---men for whom the division between the animal and the human
does not exist, men who have never wallowed in the mire, have no bad
conscience, and are in need of no conversion. These authors now come
to them and treat them like professional drunkards: For heaven's sake,
they call out to their readers, not a drop of wine to your lips! Look
around you! You live in a world of drunkards and deliriants. See,
there lies one in the gutter, dead drunk to the point of inhumanity;
there lies one trembling with twitching muscle-fibres, mortally ill in
a padded cell. For heaven's sake, empty no glass of wine! And
meanwhile grapes ripen under the sun's rays; they hang there, dark
and pale, and Dionysus---the god of wine and tragedy---stands with
green leaves around his brow and holds out to the thirsty and weary
the drink that is humanity's consolation. These moralists make Bacchus
himself into a candidate for delirium tremens. And they forget that
their words not infrequently turn to men who feel like Greeks, who
worship the wine-god, drink the wine mixed with water, enjoy the mild
intoxication, and shun drunkenness as a disgust and a plague.

Yet Tolstoy is not satisfied with the Bjørnsonesque standpoint. He
proceeds much more thoroughly. Out of his book beats a mighty
passionate fervour of feeling like flames; and his consequential
thinking goes so much deeper than Bjørnson's preaching as early
Christianity is deeper than the seminary wisdom of the doctrine of
development.

Tolstoy is offended by the animal in the human being, by the beast in
man. He is offended by the marital life of spouses. The passionate
cohabitation of the newly married destroys, in his conception, all
deeper goodwill and all genuine respect between man and woman. It is
to blame for all the evil that later surfaces between them. It produces
self-contempt and contempt for the other party. It awakens, in the
intervals between the periods of physical attraction, the deep hatred
that in his view generally breaks out between spouses in every quarrel
about even the most trifling matter, and it is moreover to blame for
the half-mad jealousy that, on the book's presupposition, always
consumes one or both parties and makes their domestic life an eternal
hell.

It is amusing in our day to encounter such an early-Christian manner
of feeling. Not a trace of the healthy, antique and modern outlook on
life, which in the cohabitation of two lovers---however passionate it
may be---sees something pure and beautiful. Not a trace of the nature
worship that in antiquity ennobled sensual life; not a trace either
of the infinite tenderness that in modern people ennobles it and
transforms crude pleasure into joy or happiness.

The starting point is the biblical word that whoever looks at a woman
with desire has already committed adultery with her in his heart. And
what is new is that for Tolstoy this word holds also within marriage.
Marriage is broken by the spouses' desire for each other. In the
perfect relationship, the woman is for the man a sister, nothing more.
The clear and consistent consequence would be the one drawn by the
religious sectarians the Skoptsy in Russia, namely self-mutilation.
Tolstoy evidently shrinks from this. But in compensation he wishes,
as far as one can understand, the sensual element either entirely
mortified or at least suppressed to the least possible degree. Not
only every exposure of the body's forms by low-cut or tight-fitting
garments, not merely dancing, is to him a horror as sensual
stimulation; he also fears art altogether, and music in particular,
less as an appeal to the senses than as an arouser of feeling. And
it is indeed after a piece of music that the book is named. The line
of thought is that every significant composition powerfully sweeps
the listener---as if by possession---into the agitated mood in which
the composer found himself when he created it; the souls of listener
and composer flow together. And thus music also merges listeners,
sometimes two by two, into this same mood, separates them from their
surroundings, and unites them. It can, in this respect, perform the
work of a procurer.

\section*{II}

The feeling of the sexual relation itself as a degradation is, after
Europe's long education in a Christian-spiritual direction, no rare
occurrence. One encounters it in our day, however, as fully developed
in personalities of a wholly pagan type as in spiritualists and
Christians. Modern development has led refined people of the most
varied convictions equally far from the Greek unity of feeling in its
freedom and health. I recall from my early youth a long philosophical
conversation with a young French friend that revolved around this very
point---sexual union---which revolted him. He was indignant that the
world's architect had not devised a more refined and noble form for
the merging of the sexes. An exclamation of his still rings in my ear:
\textit{Pourquoi cette insulte!} Why this affront, this slap in the
face!

It is this same thought that is varied in several places in Maupassant's
latest book \textit{L'inutile beauté}. In \textit{Un cas de divorce} the
hero comes, in relation to his young wife, finally to the point where
he cannot touch her with his hand or lips without feeling an
unconquerable disgust---not exactly at her, but a more comprehensive
and more contemptuous disgust at the erotic embrace itself, which for
all refined beings stands as a thing one must be ashamed of and conceal,
yes, which one mentions only in a low voice, with a blush.

And when, in the story after which the collection is named, the
conversation turns to a beautiful, refined lady whose beauty has
suffered from numerous childbirths, and the remark is made, ``What
can one say---it is nature's way!''---the young man whom the author
uses as his mouthpiece exclaims: ``Nature! I tell you that nature is
our enemy, that we must always fight against nature; it perpetually
leads us back to the animal.'' And he now develops the argument that
everything clean, beautiful, choice, and ideal on earth is not God's
but humanity's work:

\begin{quote}
``It is we who, by singing creation, interpreting it, admiring it as
poets, ennobling it as artists, explaining it as scientists, have
introduced a little grace and beauty, something attractive and
mysterious into it. Think of reproduction! What can be imagined more
ignoble and repellent than this filthy and ridiculous act, over which
all finer souls will feel revolted to the end of days. If the organs
had necessarily to serve a double purpose, why could the Creator not
have chosen, for reproduction---the noblest and most important of all
human activities---organs that were less unclean, less soiled? It is
as if, in a kind of malicious cynicism, he wished to place
insurmountable obstacles in man's way, preventing him from ever being
able to beautify or ennoble his relationship with woman. Man has
nevertheless invented love, which is not a bad reply to a cunning and
teasing deity.''
\end{quote}

Here in Maupassant, then, nature is conceived in this relation as the
enemy. In Tolstoy it is the nearest goal---though a goal beyond which
he in turn strives; an ideal lying behind us, to which one should
return, except that he knows a still higher ideal than it, namely its
renunciation. In him the struggle against sexual life is only an
application of his general perspective, according to which culture is
of the evil: money, magnificent buildings, refinement, abundance, all
life other than peasant and artisan life is of the evil. He had long
since disdained and condemned science. Now it is art's turn. He had
never been willing to see in money condensed labour, only condensed
force. Now he sees in sexual cohabitation only an animal and degrading
condition. He has brooded over this tragic and comic fact that the human
being is by its nature a more highly developed animal, so that it can
only partially consign its origin to oblivion and can never entirely
efface its traces. And he sees our culture as the power that gives
everything animal the upper hand. One feels in him, with the years,
a mounting revulsion for the merely natural, the universally human---
both Russian religiosity and the old man's gaze at these mysteries,
which has made itself felt in the North as well. The disgust at
licentiousness, at the heightened life of instinct, at passion and
vice, has in him taken the final form of loathing for the animal in
the human being.

If one passes through the book to its author and from him to the great
public that loses itself in admiration for him, the following remark
seems to me to suggest itself. It is astonishing how humanity even
today is struck with reverence for teachers of duty. It lies so deeply
in the blood of most of us, from biblical upbringing, that we
involuntarily think: here, where morality is preached, something
especially serious and solemn is occurring.

The Russians had a refined, amiable, warm-hearted author in Turgenev,
a man of the world who unpretentiously wrote one poetic masterpiece
after another. He did not imagine that he was thereby transforming the
human race. He was a quiet artist with certain weaknesses and great
virtues---a gentleman.

Then there comes after him a writer like Tolstoy. In his youth he is
a rather dissolute landowner and officer; in his more mature years he
becomes a great admonishing poet; with the years, an ever more
passionate teacher of duty; in the end, a John the Baptist. He
fashions his own religion, he preaches, he calls and rebukes, he goes
about in a camel-hair shirt---or rather in a peasant smock. He lets
his beard grow and parts his hair in the middle of his forehead like
a Russian peasant; he lives on his country estate, immersed in
contemplation and manual work; at times he resembles a peasant
prophet; and when the seriousness of his daily occupations has at
last left its deep impress on his features, he allows a visiting
photographer to fix this expression on the prepared plate. He walks
behind the plough with his white nag before the old-fashioned wooden
plough, dreams himself back into the state of nature, far from
civilization, from his manor house that lies right in the vicinity,
and allows Repin to paint him thus set back into the state of nature.

The difference is plain enough between him and the original John the
Baptist. The latter, as is well known, had no manor house, so he had
considerably less difficulty in feeling himself entirely detached from
civilization. And when he had put on the camel-hair shirt, he did not
allow himself to be photographed in it.

\section*{III}

At his highest point Tolstoy now attacks his contemporaries by virtue
of conditions whose ground lies in the human being being an animal
developed toward humanity. Could he, he would gladly eradicate the
animal in the human being entirely.

It is remarkable, this designation: the Animal. It is of course an
ancient Christian expression. One does not easily forget the vivid
depiction of the Beast in the thirteenth chapter of the Revelation of
John: ``I stood on the shore of the sea and saw a beast rise up out
of the sea. It had ten horns and seven heads; on its horns were ten
diadems and on its heads blasphemous names, and the beast I saw was
like a leopard; its feet were like a bear's, and its maw like a lion's
maw, and the dragon gave it its power, and they worshipped the Beast
and said: who is like the Beast and who can fight against it?''

The Beast was, for the Christians of that time, primarily the Roman
world power---we know this with complete precision; the symbol is
transparent to us right down to the mysterious number 666 in verse 18,
which is Nero's name, Nero Caesar, written in Hebrew letters and added
up according to the numerical value of those letters. But more deeply
conceived, the Beast was of course for the first Christians paganism
as essence and power---everything unbaptized, everything insofar
animal in the human being, that was to be overcome. To use the modern
term: the Beast was no actual animal; it was the beast in man.

That was also how Alexandre Dumas understood it when in 1870, with a
witty application of the old biblical passage, he told of his attempts
to examine the human being in the great crucible Paris. He maintained
that the human being, even the Parisian, always contained however
small a fraction of soul---roughly as the sixtieth homeopathic dilution
still contains an atom of the original fluid. And he told how he saw
male and female bovine heads being formed and shaped by the cauldron's
vapours, when he suddenly heard a boiling sound, and up from the
cauldron there rose---not formed from its foam or steam but fashioned
from the substances themselves within it---an enormous beast with
seven heads and ten horns, and on these horns ten diadems, and on
these heads hair with a metallic lustre and the colour of alcohol. The
beast resembled a leopard, its feet were powerful as a bear's, its
maw was like a lion's maw, and the dragon gave it its strength. And
the beast was clothed in purple and scarlet and adorned with gold,
with precious stones and pearls, and in its white hands it held, as
one carries a bowl of milk, a golden vase full of all the abominations
and impurities of Babylon, Sodom, and Lesbos. From its body issued an
intoxicating vapour, through which it shone like the most beautiful of
God's angels, and through which a thousand small human beings moved,
writhing with voluptuousness and howling with pain, and disappeared
with a small pop or crack---that is, they burst, and nothing remained
of them but a drop of something liquid, a tear or a drop of blood.
But the beast was not sated. It crushed them with its feet, tore them
with its claws, chewed them with its teeth, suffocated them against
its breast, and those whom it thus suffocated were the most envied.
Its seven heads formed a garland reaching to heaven; its seven mouths
were perpetually smiling and the lips flame-red, and above its ten
diadems there blazed in a blaze of light the single word: Prostitution.

This beast, he says, was the newest bodily form of woman. The Beast
here is the smiling and roaring sex that, after a thousand years of
bondage and impotence, takes revenge on the man, armed with her beauty
and her beauty's resources, and bestows on him the love that breaks
down and destroys him. And it is this Beast that Dumas threw onto the
boards of the stage when he erected the great, coarse symbol Césarine
in \textit{La femme de Claude}. For Dumas it is then, roughly as for
Tolstoy, culture---more precisely, over-culture---that makes the human
being animal in the biblical sense, that is, vicious; and the human
being was not that, on the presupposition, so long as it was not a
modern human being but subject to the unconditional discipline of
unearthly authorities.

\section*{IV}

At the same time that Tolstoy in Russia was composing \textit{The
Kreutzer Sonata}, Zola in France was writing his latest novel
\textit{La Bête humaine}---\textit{The Beast in Man}. Here the
relationship is seen from an entirely different side. For Zola the
animal element is not something that civilization develops in us, but
something that has been left behind in us as a slumbering remnant from
primordial times and that can, under certain circumstances, awaken to
our horror.

Zola's standpoint is not the moral but the natural-historical. He
presents the relationship as it appears to him, without accusation,
but with a grand-scale seriousness.

At bottom it is always the beast in man that has preoccupied Zola.
Certain younger writers like Bourget and Rod interest themselves only
in the human being in its latest and highest unfolding under a refined
civilization; Zola always reaches back to the original instincts and
sensations of which primordial humanity consists. One might think, for
example, of his depiction of the peasant's love of the land, of
property, of gain in \textit{La Terre}, or of the drive of
self-preservation in its most naive form in \textit{Le ventre de
Paris}.

Here he has reached back to one of the most sinister of original
impulses: the drive to annihilation, the cruel instinct that leads to
destruction and murder, and has conceived of this impulse in its
mysterious connection with the erotic drive. Whereas Tolstoy in his
story studies principally the sexual inclination in its over-development
and shows how the murderous madness germinates from it---the drive to
annihilation arising as a mere expression of jealousy---Zola immerses
himself principally in the murder instinct and depicts a whole series
of murders and murder attempts, but shows in all of them, as a
secondary thread, the connection with a low-born but vital eroticism.

We see the original inclination to kill alive in all the human beings
he presents. A husband who learns that an old libertine took possession
of his wife when she was a young girl avenges himself, after a fit of
retrospective jealousy, by a murder committed in a railway carriage,
and forces his wife to assist in the murder. The bloody deed lies
oppressively on the couple and soon drives them entirely apart.

The young wife takes a lover and demands that he kill her husband. He
consents knowingly, but at the very moment he lies in wait for the
husband's arrival in order to murder him, he kills, driven by an
irresistible impulse and against his will, his mistress.

This young man, the involuntary murderer, is the main character of
the book. Zola has, as he has expressly stated, sought in contrast to
the Russians Dostoevsky and Tolstoy to show that one murders not with
deliberation and clear consciousness but from instinct. He wanted to
furnish a kind of counterpart to Raskolnikov, and comes remarkably
close to Tolstoy in doing so.

Jacques Lantier suffers from a morbid drive, a mysterious and
irresistible impulse that is found in nature, is well known to
physicians, and has frequently been described in learned handbooks.
But what Zola has added out of his own as a personal peculiarity, as
an improbable explanation, is this: what comes over Lantier is not
simply a sudden seizure of blind rage but a continually recurring
thirst to avenge very old wrongs of which he has lost all precise
memory. It came from so far back, the text has it---from all the evil
women had done to his forebears, from the grievance that had
accumulated from one male to another, ever since a woman betrayed one
of them for the first time, when they still inhabited caves. The fits
are thus regarded entirely as a haunting; they stem from prehistoric
times, quite opposite to Pozdnyshev's fits in Tolstoy, which are
derived from the over-civilization of modern times. When Lantier has
his fits, he feels it as a necessity to struggle to conquer a woman
and possess her entirely. This is, for Zola's imagination, the
degenerate impulse to throw her over his back as prey---the impulse
the male has always wrested for himself from community with the others.

In Zola, then, imperfectly digested Darwinism---as in Tolstoy,
irrationally revived early Christianity.

Grouped around Lantier are other characters with similar potential for
murder: a young woman, a railway gatekeeper, who out of jealousy
causes an entire railway train to be wrecked because she is in love
with Lantier and wants to revenge herself on him and his mistress, who
are both in the train; an old man, a moral monster, who slowly
poisons his wife in order to get hold of her savings; and finally a
fireman who out of jealousy at last murders Jacques.

This, then, is a great poem about prehistoric murder-lust in its
connection with the erotic drive---in that respect a kind of
prehistoric epic. It is instructive to compare it with a genuinely
prehistoric epic of our own day, such as the German Heinrich Hart's
\textit{Tul und Nahila}, the first volume of a work entitled \textit{Lied
der Menschheit}, which after its naive plan is to consist of no fewer
than twenty-four cantos.

Hart's poem takes place in Ceylon in the grey antiquity of thousands
of years before all civilization. The main characters are a human
couple living among beings from the transitional time between
humanity's ancestors and humans themselves. Ape-like dwarfs and giants
still carry on their eerie game alongside the humans. Only a few
concepts have developed: fire is the fire-serpent; how to master it
is still unknown. Expressions are used like ``the fire-serpent has
gnawed at the flesh.'' Few tools: club, stone axe, stone-tipped spear,
bone fish-hook, bast sling for hunting. Among the humans, natural
dependence and herd-instinct hold sway. We follow the life-history of
the human couple. Nahila loves Tul for life because this man, for her
sake, was expelled from his tribe when he refused to yield her to the
community. A sense of shame, of modesty, develops in her from a
peculiar and not improbable motive: because her own sex seems ugly to
her---therefore she makes herself a skirt of leaves. We see how the
birth of a child makes the taming of domestic animals necessary under
certain conditions. A goat is caught that nourishes the child with its
milk. Despite all the difficulties the subject presents, the poem is
not stupid. The human beings seem here genuinely original and wild.
And yet everything in Zola is, in its kind, wilder and greater, even
though it takes place in our own day---in 1870---and appears as a
novel about railway operations and railway accidents and about crimes
committed with railway tracks or in railway carriages.

Jules Lemaître has pointed out the weakness in this book: that the
setting has no inherent connection with the core. Whereas in
\textit{L'Assommoir} there is a necessary connection between the
dram-shop and the working-class population, or in \textit{Germinal}
between the mine and the miners, here there is no necessary connection
between the lust for annihilation and railway operations. To that
extent this book is inferior to several preceding ones, more arbitrary
in its composition. But it is great through the author's eye for the
simple and terrible basic drives of humanity, and great through its
rich symbolism.

Besides the beast in man, there is in this novel, as in Zola's others,
an impersonal being that appropriates a major share of the author's
interest. In earlier books there appeared the cemetery, the market
halls, the dram-shop, the department store, the farmland, the church.
We have seen above that while Zola consistently traces the human back
to the animal, he attributes to impersonal creatures all the more-than-
animal force of which he deprives the individual beings---indeed,
attributes to them independence and will. Such an impersonal creature
in the railway novel is a locomotive. It has a name; it is called
Lison; it leads a life and arouses feelings like a woman. The engine
driver loves his machine, loves it as the rider loves his horse. A
kind of marriage takes place between them.

Here there occurs a magnificent description of the locomotive's
struggle to push through accumulated snow in a storm, and when the
train later collides with a cart loaded with heavy stones and the
locomotive is wrecked, it dies just like a living being. There is in
this an art that has become artfulness. The locomotive, this monster
with glowing eyes and iron flanks, puts the human beings to shame, so
to speak, by its docility, its masculine discipline, its whole good
behaviour, and by the unmixed love it inspires in its driver.

And at the end of the book there is no lack of another, greater symbol.
When the fireman has murdered Jacques, pushed him down from his
locomotive, and when this---without a driver, without a guiding
spirit---races like a mad creature past station after station, carrying
French conscript soldiers to the border in the July of that remarkable
year 1870, it stands written:

\begin{quote}
``It rolled, rolled on through the black night; nobody knew where. What
did it matter how many victims the machine crushed on the way? Did it
not roll on, in spite of all, toward the future, heedless of the blood
spilled? Without a driver, in the darkness, a blind and deaf animal,
it rolled and rolled on, laden with this cannon-fodder, these already
exhausted soldiers who sang in their drunkenness.''
\end{quote}

This locomotive that is guided by no one is, as one feels, the symbol
of that France that goes to war without a will at its head. The novel
about the drive to annihilation, about the murder instinct, reaches
its culmination in this dragging of thousands upon thousands toward
this enormous mass murder. The endless bloodshed of the great war is
the triumph and defeat of the beast in man in one.

\section*{V}

What is admirable about these great foreign authors and writers is
that they unfold the human being before our eyes completely, like a
fan. There is nothing that, out of fear of readers or reviewers, needs
to be concealed or falsified. We look out over a wide field of
spiritual life, from the vicious and low, from crime, through the
animally human, the merely natural, the universally human, to the
purely human and its human nobility. And everything that is communicated
is communicated in the admirable freedom of expression of great
literatures, without prudishness and without sentimentality.

All of them need, in order to give an impression of life's richness
and misery, the idea of a split in our nature into two worlds---the
unconscious and the conscious---which in the poetic language of most
recent times have been made synonymous with the worlds of the animal
and the human within us, the domains of necessity and freedom, as it
was said in former times.

Modern science no longer believes in a freedom of will of the kind
previously clung to. Some of the writers see a danger precisely in
this. Thus Bourget, who seems frightened by the impression that the
philosophical doctrine of necessity has sometimes made on a younger
generation, in whose eyes it provides a licence for ruthless egotism.
He fears that science itself, civilization itself, will unleash the
animal in coming generations, and in his last books (\textit{Mensonges},
\textit{Le Disciple}) has, with a certain pietistic flaccidity,
referred to the Catholic abbé, to religion, to prayer, as the only
power capable of holding the animal down and giving modern minds unity
anew.

Dumas, who is enough of a genuine Frenchman to harbour no doubt about
the freedom of the will, sees the animal develop in the swamp of over-
civilization and cries ``guard!'' against it. The moralist Tolstoy and
the child of the world Maupassant regard with the same loathing the
fundamental relation between the two sexes, which for them both stands
as a violation of human dignity---but which for Maupassant is
degrading as a piece of unconquered and unconquerable primordial nature,
and for Tolstoy degrading as a product of culture's over-stimulated
unnaturalness. Zola, finally, is the thoroughly convinced believer in
necessity, for whom the human plant, the beast in man, human thought,
human will, obey unshakeable laws, and who finds in compensation a
kind of fantastic satisfaction in equipping lifeless objects with
independent life and autonomous will.

Human, truly human, is only the machine, Zola seems to say in his
latest novel. Human, truly human, is only civilization, Maupassant
seems to say in his latest collection of stories, which passes judgment
against everything that is nature in us, as animal. Human, nobly human
---in the language of former times, divine---is only the soul's
recollection in devout feeling and prayer, says Bourget in his two
latest books. Human is not culture, not civilization, says Tolstoy in
his latest works, but only the way of life that never meets evil with
self-defence, never wards off injustice, never troubles itself about
property and money, and that brings man and woman always to see in
each other only a sister and brother.

Such, then, is the language spoken by the wise of our day, by the men
who in our time address the greatest and most distinguished readership
on earth.

If one could imagine that an ancient Hellene from the classical period
of Greece were to rise from his ashes and hear these pronouncements,
he would surely clap his hands together in astonishment at such talk.

\end{document}
